This section describes the statistical framework used to infer the network structure underlying the observed data. The approach is based on Gaussian graphical models and Bayesian inference over the space of decomposable graphs. This section is a direct review of the method proposed in \textcite{donnet2012empirical}, for a more in-depth methodological explanation, refer to the original paper or to  \textcite{lauritzen1996graphical} for a textbook treatment.
\subsection{Primer on Decomposable Gaussian Graphical Models}
Let $G = (V,E)$ be an undirected graph with vertices $V = \{1,...,p\}$ and $E=\{e_1,...,e_t\}$. A graph or subgraph is said to be complete if all pairs of its vertices are joined by edges. A complete subgraph that is not contained within another complete subgraph is called a clique. Let $\mathcal{C}=\{C_1,...,C_k\}$ be the set of cliques of an undirected graph. An order of the cliques is said to be perfect if $\forall i =2,...k$, $\exists h = h(i) \in \{1,...,i-1\}$ such that $S_i = C_i \cap \cup_{j=1}^{i-1} C_j \subseteq C_h$. $\mathcal{S}=\{S_2,\dots,S_k\}$ is the set of separators associated to the perfect order. An undirected graph admitting a perfect order is said to be decomposable (or equivalently chordal). See Figure~\ref{fig:decomposability_cliques} for a visual example of a decomposable graph\\

\begin{figure}[h]
    \centering
    % Define layers so we can draw colored blobs behind the nodes
    \pgfdeclarelayer{background}
    \pgfsetlayers{background,main}

    \begin{tikzpicture}[scale=1.6, auto, swap]
        % =========================================
        % LEFT: Non-Decomposable (The "Hole")
        % =========================================
        \begin{scope}[xshift=-2.5cm]
            % Nodes
            \node[draw, circle, fill=gray!20, thick] (1) at (0,2) {1};
            \node[draw, circle, fill=gray!20, thick] (2) at (2,2) {2};
            \node[draw, circle, fill=gray!20, thick] (3) at (2,0) {3};
            \node[draw, circle, fill=gray!20, thick] (4) at (0,0) {4};

            % Edges forming the 4-cycle
            \draw[thick, red] (1) -- (2) -- (3) -- (4) -- (1);

            % Label
            \node[align=center, below] at (1, -0.5) {\textbf{(a) Non-Decomposable Graph}\\ \footnotesize(A ``hole'': chordless 4-cycle)};

            % Visual cue for the hole
            \node[red, font=\huge] at (1,1) {$\times$};
        \end{scope}


        % =========================================
        % RIGHT: Decomposable with Cliques & Separators
        % =========================================
        \begin{scope}[xshift=2.5cm]
            % 1. Draw the Cliques in the background layer

            % 2. Draw Nodes on top (same positions as left)
            \node[draw, circle, fill=white, thick] (1b) at (0,2) {1};
            \node[draw, circle, fill=white, thick] (2b) at (2,2) {2};
            \node[draw, circle, fill=white, thick] (3b) at (2,0) {3};
            \node[draw, circle, fill=white, thick] (4b) at (0,0) {4};


                        \begin{pgfonlayer}{background}
                % Clique C1 = {1, 2, 3} (Top-Right Triangle) - Shaded Blue
                % Using a smooth plot to draw a blob around the nodes
                \filldraw[blue!30, opacity=0.5, draw=blue, dashed, line width=1pt]
                    plot [smooth cycle, tension=0.2] coordinates {
                        ($(1b)+(-0.2, 0.3)$)  % Above node 1
                        ($(2b)+(0.3, 0.3)$) % Top-right of node 2
                        ($(3b)+(0.3, -0.3)$) % Bottom-right of node 3
                    };

                % Clique C2 = {1, 3, 4} (Bottom-Left Triangle) - Shaded Red
                \filldraw[red!30, opacity=0.5, draw=red, dashed, line width=1pt]
                    plot [smooth cycle, tension=0.2] coordinates {
                        ($(1b)+(-0.2, 0.3)$) % Top-left of node 1
                        ($(3b)+(0.1, -0.2)$) % Bottom-left of node 3
                        ($(4b)+(0, -0.2)$)  % Below node 4
                    };
            \end{pgfonlayer}

            


            
            % 3. Draw Edges
            % Standard edges
            \draw[thick] (1b) -- (2b) -- (3b) -- (4b) -- (1b);
            % The Chord that makes it decomposable (highlighted)
            \draw[line width=2pt, purple!80!black] (1b) -- (3b);

            % 4. Add Labels for Cliques and Separator
            % Label C1 near node 2
            \node[blue!80!black, right] at ($(2b)+(0.4,0)$) {\large $C_1=\{1,2,3\}$};

            % Label C2 near node 4
            \node[red!80!black, left] at ($(4b)+(-0.4,0)$) {\large $C_2=\{1,3,4\}$};

            % Label Separator pointing to the chord
            % The separator is the intersection: {1,3}
            \node[purple!80!black, align=center] (sepLabel) at (3, 1) {Separator\\$S_1 = C_1 \cap C_2$\\$S_1=\{1,3\}$};
            \draw[-{[length=3mm]}, purple!80!black, thick, bend right=20] (sepLabel.west) to ($(1b)!0.5!(3b)$);


            % Main Subfigure Label
            \node[align=center, below] at (1, -0.5) {\textbf{(b) Decomposable Graph}\\ \footnotesize(Decomposed into cliques)};
        \end{scope}

    \end{tikzpicture}
    \caption{Visualizing Decomposability. Graph (a) contains a forbidden ``hole''. Graph (b) is separable by adding the chord (1,3). This allows the graph to be decomposed into two maximal cliques, $C_1$ (blue shaded area) and $C_2$ (red shaded area). The separator $S_1$ (purple highlighted edge) is the intersection $C_1 \cap C_2$, acting as the bridge that conditionally causes independence between nodes 2 and 4.}
    \label{fig:decomposability_cliques}
\end{figure}

To each vertex $v \in V$ we associate a random variable $y_v$. Thus for $A \subseteq V$, $y_A$ denotes the collection of random variables $\{y_v : v\in A\}$. To simplify, let $y = y_V$. The probability distribution of $y$ is said to be Markov with respect to $G$ if for any decomposition $(A,B)$ of $G$, $y_A$ is independent of $y_B$ given $y_{A\cap B}$. A graphical model is defined as a family of distributions on $y$ verifying the Markov property with respect to a graph.\\
A Gaussian graphical model is such that:
\begin{equation*}
    y|G, \Sigma_G \sim \mathcal{N}_p(\mu, \Sigma_G)
\end{equation*}
Where $\mu \in \mathbb{R}^p$ and $\Sigma_G$ is the associated $p\times p$ symmetric positive definite covariance matrix. In the Gaussian case, $y$ is Markov with respect to $G$ if and  only if 
\begin{equation*}
    (i,j) \not\in E \iff (\Sigma_G^{-1})_{(i,j)} =0
\end{equation*}
The matrix $\Sigma_G^{-1}$ is referred to as the concentration or precision matrix. The density of $Y = (y^1,...,y^n)$ from the previously defined model when demeaned is a function of multivariate Gaussian densities on the cliques and separators of $G$. Let $\mathcal{C}$ and $\mathcal{S}$ denote the sets of cliques and separators of $G$ respectively, corresponding to a perfect order for $G$. Then:
\begin{equation*}
    \begin{gathered}
        f(Y|\Sigma_G,G) = \prod_{i=1}^n \{\frac{\prod_{c\in\mathcal{C} }\phi_{|C|}(y_C^i | (\Sigma_G)_C)}{\prod_{S\in \mathcal{S}}\phi_{|S|} (y_S^i | (\Sigma_G)_S)}\}
    \end{gathered}
\end{equation*}
Such that $(\Sigma_G)_A$ is the submatrix corresponding to $\{(\Sigma_G)_{i,j}\}_{i\in A, j \in A}$ and $\phi_q(\cdot| \Delta)$ is the q-variate Gaussian density with mean $\mathbf{0}_q$ and $q\times q$ covariance matrix $\Delta$. In the Bayesian framework, we are interested in the graph posterior probabilities of the form:
\begin{equation*}
    \pi(G | Y) \propto \pi (G) \int f(Y|\Sigma_G , G) \pi(\Sigma_G|G)d\Sigma_G
\end{equation*}

\subsection{Prior specification}
For the prior on the covariance matrix, conditional on $G$, we set a Hyper-Inverse Wishart (HIW) distribution prior on $\Sigma_G$:
\begin{equation*}
    \Sigma_G| G,\delta, \Phi \sim HIW_G (\delta, \Phi)
\end{equation*}
This distribution entails (See \textcite{lauritzen1996graphical}) that for every clique $C \in \mathcal{C}$, $(\Sigma_G)_C \sim IW(\delta, \Phi_C)$ with density:
\begin{equation*}
    \pi((\Sigma_G)_C | \delta, \Phi_C) = h_{G_C}^{IW} [\det (\Sigma_G)_C]^{-\frac{\delta+2|C|}{2}}\exp\{-\frac{1}{2}\text{tr}[(\Sigma_G)_C^{-1} \Phi_C]\}
\end{equation*}
Where $h_{G_C}^{IW}$ is the following normalization constant:
\begin{equation*}
    h_{G_C}^{IW}(\delta, \Phi_C) = \frac{\det(\frac{\Phi_C}{2})^{(|C|+\delta-1)/2}}{\Gamma_{|C|}(\frac{|C| +\delta - 1}{2})}
\end{equation*}
Where $\Gamma_v$ is the multivariate gamma function\footnote{See \href{https://en.wikipedia.org/wiki/Multivariate_gamma_function}{Multivariate Gamma Function} for more details.} with parameter $v$. Conditional on $G$, the HIW distribution is conjugate and the posterior distribution of $\Sigma_G$ is given by $\Sigma_G | Y,G, \delta, \Phi \sim HIW(\delta+n, \Phi + S_Y)$. Where $S_Y$ is the empirical covariance matrix. For such a prior, the marginal likelihood for $G$ turns out to be a simple function of the HIW prior and posterior normalizing constants:
\begin{equation}\label{margData}
    f(Y|G, \delta,\Phi) = \frac{h_G(\delta, \Phi)}{(2\pi)^{-np/2}h_G (\delta+n, \Phi + S_Y)}
\end{equation}
Where the normalizing constant $h_G$ can be explicitly computed in decomposable graphs from the normalizing constant for the cliques and separators:
\begin{equation*}
    h_G(\delta,\Phi) = \frac{\prod_{C\in \mathcal{C}}h_{G _C}^{IW} (\delta, \Phi_C)}{\prod_{S\in \mathcal{S}}h_{G_S}^{IW} (\delta, \Phi_S)}
\end{equation*}
For the graph itself, the prior of choice is a Bernoulli distribution, with parameter $r$ on the inclusion or not of each edge:
\begin{equation}\label{priorG}
    \pi(G|r) \propto r^{k_G} (1-r)^{m-k_G}
\end{equation}
Where $k_G$ is the number of edges on $G$ and $m = \frac{p(p-1)}{2}$ is the number of possible edges.
Using (\ref{priorG}) and (\ref{margData}) it can be shown that the density of the posterior distribution in the space of decomposable graphs satisfies:
\begin{equation}\label{postGraph}
    \pi(G|Y,\delta,r,\Phi) \propto \frac{h_G (\delta,\Phi)}{h_G(\delta+n, \Phi + S_Y)}\pi(G|r)
\end{equation}
\subsection{Choice of hyper-parameters $\delta,r$ and $\Phi$}
Following \textcite{donnet2012empirical}, we set $\delta = 1$, which is equivalent to a prior weight equal to one observation. For the structure of $\Phi$, we set $\Phi = \tau I_p$, where $\tau$ is a shrinkage factor that controls the sparsity of the implied graph. The choice of $\tau$ is therefore crucial and must be calibrated to balance sparsity against structure. We adopt an empirical Bayes strategy that fixes $(\tau, r)$ at their maximum likelihood estimates. Because direct maximization is computationally challenging, the update is incorporated into the EM step of the stochastic approximation EM (SAEM) algorithm that drives the SAEM-MCMC procedure. The SAEM algorithm, described in detail in the next section, is a stochastic approximation to a standard EM algorithm.